%% kapa: exam=1 round=34 year=2023 period=1 form=A subject=real_estate subject_ko=부동산학원론 seq=12 number=92
%% source: Q-Net 감정평가사 기출문제 (cst00309) artlSeq=5211789 file=raw/1008/5211789/2023년도 제34회 감정평가사 1차시험 1교시.pdf
%% conversion: kordoc 4.24.51 -> tools/gen1.py (auto); review=proofread-sonnet; flags=none
\begin{question}{92}
A도시와 B도시 사이에 있는 C도시는 A도시로부터 5km, B도시로부터 10km 떨어져 있다. 각 도시의 인구 변화가 다음과 같을 때, 작년에 비해 금년에 C도시로부터 B도시의 구매활동에 유인되는 인구수의 증가는? (단, 레일리(W. Reilly)의 소매인력법칙에 따르고, C도시의 모든 인구는 A도시와 B도시에서만 구매하며, 다른 조건은 동일함)

{\small\centering\renewcommand{\arraystretch}{1.4}
\begin{tabular}{|c|c|c|}\hline
구분 & 작년 인구수 & 금년 인구수 \\ \hline
A 도시 & 5만 명 & 5만 명 \\ \hline
B 도시 & 20만 명 & 30만 명 \\ \hline
C 도시 & 2만 명 & 3만 명 \\ \hline
\end{tabular}\par}
\begin{choices}
  \item 6,000명
  \item 7,000명
  \item 8,000명
  \item 9,000명
  \item 10,000명
\end{choices}
\end{question}
\answer{3}
